% ============================================================
%  chapter/chapter2.tex — Figures & tables stress test
%
%  Real artwork goes into graphs/ and illustrations/.  Until then
%  the \placeholder boxes show where images will sit.
% ============================================================
\chapter{Figures and Tables}\label{ch:figures}

\section{Figures}
Drop your plots into \texttt{graphs/} and load them with
\texttt{\textbackslash includegraphics}; the placeholder below shows the frame the
finished figure will occupy.

\begin{figure}[ht]
  \centering
  % Real plot: \includegraphics[width=0.8\linewidth]{graphs/plot1.jpg}
  \placeholder[AccentColor!12]{0.8\linewidth}{55mm}
  \caption{A placeholder for a data plot. Replace the
    \texttt{\textbackslash placeholder} line with an
    \texttt{\textbackslash includegraphics} pointing at
    \texttt{graphs/plot1.jpg}.}
  \label{fig:plot1}
\end{figure}

Figure~\ref{fig:plot1} floats to the top of a page when that helps
the layout; \texttt{[ht]} lets LaTeX choose between ``here'' and
``top''. Side-by-side sub-figures work too:

\begin{figure}[ht]
  \centering
  \begin{subfigure}{0.48\linewidth}
    \centering
    \placeholder{0.9\linewidth}{35mm}
    \caption{First panel.}
  \end{subfigure}\hfill
  \begin{subfigure}{0.48\linewidth}
    \centering
    \placeholder{0.9\linewidth}{35mm}
    \caption{Second panel.}
  \end{subfigure}
  \caption{Two sub-figures sharing one number.}
  \label{fig:panels}
\end{figure}

\section{Tables}
Booktabs rules (no vertical lines) are the professional standard:

\begin{table}[ht]
  \centering
  \caption{A small comparison table.}
  \label{tab:compare}
  \begin{tabular}{lcc}
    \toprule
    Format & Trim size & Typical use\\
    \midrule
    A4       & 210 × 297 mm & review drafts\\
    B5       & 176 × 250 mm & digest editions\\
    US Trade & 152 × 229 mm & trade paperback\\
    \bottomrule
  \end{tabular}
\end{table}

Table~\ref{tab:compare} uses fixed columns; for text-heavy tables
that must fit the measure exactly, use \texttt{tabularx}:

\begin{table}[ht]
  \centering
  \caption{A tabularx table — the X column absorbs the spare width.}
  \label{tab:x}
  \begin{tabularx}{\linewidth}{lX}
    \toprule
    Term & Meaning\\
    \midrule
    Verso & The left-hand, even-numbered page of a spread.\\
    Recto & The right-hand, odd-numbered page of a spread.\\
    Folio & The page number itself, printed on the outer edge here.\\
    \bottomrule
  \end{tabularx}
\end{table}

\section{A wrapped side figure}
Narrow figures can sit beside the running text instead of floating
away from it --- useful for portraits and small diagrams:

\begin{wrapfigure}{r}{0.42\linewidth}
  \centering
  % Real portrait: \includegraphics{illustrations/portrait.jpg}
  \placeholder{0.9\linewidth}{34mm}
  \caption{A figure wrapped by the text.}
  \label{fig:wrap}
\end{wrapfigure}

Figure~\ref{fig:wrap} occupies the right of the measure while these
paragraphs flow around it. Keep wrapped figures narrow --- about a
third to two-fifths of the text width --- so the remaining lines stay
long enough to read comfortably. Two or three paragraphs is the most
text that should share a column with a wrapfigure; beyond that the
page starts to look pinched. Note how the wrapping has now finished:
a section heading should never have to share its lines with a
floating column, so leave enough prose after the figure to close the
wrap before the next heading arrives.

\section{Spanning tables and grouped rows}
Two table problems come up constantly: cells that span several rows,
and tables longer than a page. The first is \texttt{multirow}'s job,
shown in Table~\ref{tab:grouped} --- which, unlike the tables so
far, is \emph{not} a float: it is set exactly here with
\texttt{\textbackslash captionof}, so it can never drift away from
the sentence that introduces it. Use this for small tables whose
position matters; let everything else float.

% A non-floating table: \captionof (from the caption package) keeps
% the numbering and the List of Tables entry without any floating.
\begin{center}
  \captionof{table}{Rows grouped under spanning cells.}
  \label{tab:grouped}
  \begin{tabular}{llr}
    \toprule
    Family & Face & Weights\\
    \midrule
    \multirow{3}{*}{TeX Gyre} & Pagella & regular, bold, italic\\
                              & Heros   & regular, bold\\
                              & Termes  & regular, bold, italic\\
    \addlinespace
    \multirow{2}{*}{Latin Modern} & Roman & regular, bold\\
                                  & Sans  & regular, bold\\
    \bottomrule
  \end{tabular}
\end{center}

The second problem --- a table that outgrows the page --- belongs to
\texttt{longtable}, which breaks across pages and repeats its header
on every continuation:

\begin{longtable}{@{}p{28mm}p{\dimexpr\linewidth-28mm-12pt\relax}@{}}
  \caption{A long table that spans pages.}\label{tab:long}\\
  \toprule
  Term & Definition\\
  \midrule
  \endfirsthead
  \multicolumn{2}{@{}l}{\small\itshape Table~\ref{tab:long}, continued}\\
  \toprule
  Term & Definition\\
  \midrule
  \endhead
  \bottomrule
  \endfoot
  Ascender   & The part of a letter that rises above the x-height, as in b, d, k.\\
  Baseline   & The invisible line on which letters appear to sit.\\
  Descender  & The part that drops below the baseline, as in p, q, y.\\
  Kerning    & Selective adjustment of space between specific letter pairs.\\
  Ligature   & Two or more letters cast as one glyph, such as fi or fl.\\
  Serif      & The small finishing stroke at the ends of a letter's main strokes.\\
  Terminal   & The often-rounded end of a stroke, as in the tail of an a.\\
  X-height   & The height of the lowercase letters, excluding ascenders.\\
\end{longtable}
